\documentclass[12pt]{article}
\usepackage{amsmath, amssymb, amsfonts}
\usepackage{graphicx}
\usepackage{booktabs}
\usepackage{float}
\usepackage{setspace}
\usepackage{geometry}
\usepackage{url}
\geometry{margin=1in}

\title{Depression, Study Effort, and Productivity:\\ 
An Empirical and Dynamic Programming Analysis}
\author{Prachi Mallick}
\date{\today}

\begin{document}

\maketitle
\doublespacing

\section{Introduction}

Depression is known to affect motivation, concentration, and overall academic
performance. In this project, I study how depression relates to students'
study behavior and GPA outcomes using a dataset of 27,901 university students.
I first document basic differences between depressed and non-depressed students
in study hours, sleep, and financial stress. I then estimate a GPA production
function that shows depressed students have lower marginal productivity of
study time.

To better understand how these productivity differences translate into study
choices, I build a simple dynamic programming model in which students choose
daily study hours to balance GPA gains against effort costs. Using parameters
calibrated from the data, the model predicts that depressed students optimally
choose to study more hours, despite being less productive. This project combines
empirical analysis with a behavioral model to explain why effort patterns differ
across depression status.


% ============================================================
\section{Data and Cleaning}

The analysis uses the \textit{Student Depression Dataset} (Zubair Dhuddi, 2023), 
a publicly available Kaggle dataset containing self-reported information on 
demographics, study habits, mental health indicators, and academic outcomes for 
27{,}901 university students across multiple Indian cities. Because the dataset 
represents only Indian students, external validity should be interpreted within 
this geographic and institutional context.

\subsection{Key Variables}

The core variables used in estimation are:
\begin{itemize}
    \item $CGPA_i$: cumulative grade point average (0--10 scale).
    \item $Hours_i$: daily study hours (cleaned from the ``Work/Study Hours'' field).
    \item $Dep_i \in \{0,1\}$: indicator for self-reported depression.
    \item $Sleep\_clean$: numeric sleep duration (cleaned from text responses).
    \item Demographic controls: Age, Gender, Degree major.
\end{itemize}

\subsection{Cleaning Procedures}

The raw dataset contains heterogeneous text formats and several fields requiring 
normalization. All cleaning procedures are implemented programmatically in Python 
to ensure full reproducibility. The main steps are:

\begin{enumerate}
    \item \textbf{Sleep Duration}: Converted to numeric hours using lightweight 
    text parsing (e.g., midpoints of ranges, \texttt{hh:mm} conversion, numeric 
    extraction, and recognition of simple spelled-out numbers).

    \item \textbf{Study Hours}: The ``Work/Study Hours'' field is standardized 
    into a numeric variable. Ranges, \texttt{hh:mm} formats, and standalone 
    numbers all parse successfully for the full sample.

    \item \textbf{Financial Stress}: Non-numeric tokens (e.g., ``?'', ``NA'') are 
    mapped to missing values. For the remaining entries, the first numeric value 
    is extracted to construct a continuous measure.

    \item \textbf{Binary Variables}: Yes/No fields (suicidal thoughts, family 
    mental illness, depression) are standardized into binary indicators using a 
    robust parser that handles common spelling variants and punctuation.
\end{enumerate}

Quality-assurance checks confirm that:
\begin{itemize}
    \item Study Hours convert successfully for all 27{,}901 observations,
    \item 27{,}898 observations contain valid numeric Financial Stress values 
          (the remainder consist primarily of non-informative entries).
\end{itemize}

\footnote{Data sourced from Zubair Dhuddi (2023), \textit{Student Mental
Health \& Depression Indicators Data}, Kaggle. Available at:
\url{https://www.kaggle.com/datasets/zubairdhuddi/student-dataset}.}

% ============================================================
\section{Descriptive Results}

Table~\ref{tab:desc} presents summary statistics for the main variables by
depression status. Of the 27{,}901 students in the sample, 16{,}336 
(58.5\%) report experiencing depression, while 11{,}565 do not. Mean CGPA is 
similar across groups (7.62 vs.\ 7.68), with depressed students exhibiting a 
slightly higher average. More notable differences appear in study behavior and
financial well-being: depressed students report studying approximately 1.6 more
hours per day and exhibit substantially higher financial stress. Sleep duration
is lower among depressed students, consistent with patterns often associated 
with academic pressure and mental health concerns.

\begin{figure}[H]
\centering
\includegraphics[width=0.75\textwidth]{./analysis_outputs/cgpa_density_by_depression.png}
\caption{CGPA distribution by depression status (density plot). 
Depressed and non-depressed students have similar CGPA distributions; the small mean difference is visible but modest relative to within-group spread.}
\label{fig:cgpa_density}
\end{figure}


\begin{table}[H]
\centering
\begin{tabular}{lcc}
\toprule
 & Non-Depressed & Depressed \\
\midrule
Number of students & 11{,}565 & 16{,}336 \\
Mean CGPA & 7.62 & 7.68 \\
Mean study hours & 6.24 & 7.81 \\
Mean sleep hours & 6.40 & 6.18 \\
Mean financial stress & 2.52 & 3.58 \\
\bottomrule
\end{tabular}
\caption{Descriptive statistics by depression status.}
\label{tab:desc}
\end{table}

To formally assess these differences, we conduct Welch two-sample $t$-tests for
CGPA, study hours, and sleep duration. All three variables differ significantly 
across depression status. Depressed students study, on average, 
1.57 additional hours per day ($p < 0.001$), sleep 0.22 fewer hours 
($p < 0.001$), and have a small but statistically significant higher CGPA 
(difference of 0.066, $p < 0.01$). These differences suggest meaningful 
behavioral and environmental distinctions between the two groups.

Figure~\ref{fig:hours_box} visualizes the distribution of reported study hours 
by depression status. Depressed students not only study more on average but 
also display considerably greater dispersion: the upper tail of the 
distribution extends beyond 10--12 hours, while non-depressed students have a 
more compressed range. This pattern aligns with the elevated financial stress 
reported by depressed students and may reflect compensatory study behavior or 
greater academic pressure.


\begin{figure}[H]
\centering
\includegraphics[width=0.6\textwidth]{./analysis_outputs/hours_by_depression_boxplot.png}
\caption{Distribution of reported study hours by depression status. 
Depressed students report both higher average study hours and greater 
dispersion.}
\label{fig:hours_box}
\end{figure}


% ============================================================
\section{GPA Production Function}

We model academic performance as the outcome of a simple production function in 
which study hours contribute to GPA with diminishing returns. Let $H_i$ denote 
daily study hours for student $i$ and $Dep_i \in \{0,1\}$ indicate depression. 
The GPA production function is specified as:

\[
CGPA_i = \beta_0 
        + \beta_1 H_i 
        + \beta_2 H_i^2
        + \gamma_1 Dep_i
        + \gamma_2 (Dep_i \cdot H_i)
        + \gamma_3 (Dep_i \cdot H_i^2)
        + u_i,
\]

where $\beta_1 > 0$ and $\beta_2 < 0$ capture diminishing marginal returns to 
study effort. The interaction terms allow both the level and slope of the 
production function to differ for depressed students.

\begin{figure}[H]
\centering
\includegraphics[width=0.8\textwidth]{./analysis_outputs/hours_vs_cgpa_scatter.png}
\caption{Scatter plot of study hours and CGPA with group-specific fitted lines.
The fitted line for depressed students is flatter, indicating lower marginal returns to study hours.}
\label{fig:hours_scatter}
\end{figure}


% ============================================================
\section{Marginal Productivity of Study Time}

The marginal product of an additional hour of study is:

\[
MP_i =
\frac{\partial CGPA_i}{\partial H_i} =
\begin{cases}
\beta_1 + 2 \beta_2 H_i, 
    & \text{if } Dep_i = 0, \\[6pt]
(\beta_1 + \gamma_2) 
    + 2(\beta_2 + \gamma_3) H_i, 
    & \text{if } Dep_i = 1.
\end{cases}
\]

Thus, depression changes both the overall GPA level and the effectiveness of additional study hours. These marginal products are used later in the dynamic 
programming analysis to determine optimal study-hour choices

Using the estimated coefficients from the interaction regression, the marginal 
products evaluated at the mean study hours of each group are:

\[
MP_{\text{non-dep}}(H=6.238) = 0.003239,
\qquad
MP_{\text{dep}}(H=7.808) = 0.002106.
\]


A convenient measure of the productivity gap is the ratio
\[
e \;=\; 
\frac{MP_{\text{dep}}}{MP_{\text{non-dep}}}
\;\approx\;
\frac{0.002106}{0.003239}
\;\approx\; 0.65.
\]

\begin{figure}[H]
\centering
\includegraphics[width=0.8\textwidth]{./analysis_outputs/marginal_product_plot.png}
\caption{Marginal product of study hours for depressed and non-depressed students. 
The non-depressed marginal product declines with hours, whereas the depressed marginal product starts lower but increases with hours in the estimated quadratic specification.}
\label{fig:mp_plot}
\end{figure}


Thus, depressed students operate at roughly 65\% of the hourly productivity of
non-depressed students. This reduction in marginal product plays a central role
in the dynamic programming analysis, as it shifts the trade-off between effort
and GPA gains and alters the optimal study-time policy.

% ============================================================
\section{Economic Interpretation}

The marginal product calculations indicate that an additional hour of study 
raises GPA only slightly for both groups, but the gain is systematically lower 
for depressed students. At their respective average study levels, 
non-depressed students achieve a marginal product of $0.00324$, while depressed 
students achieve only $0.00211$. In economic terms, depression shifts the 
productivity schedule downward: each additional hour of effort yields smaller 
returns. This implies that depressed students face a steeper trade-off between 
effort and performance, since they must exert more hours to achieve the same 
increment in GPA. These diminished returns help explain the higher study hours 
observed among depressed students and motivate the dynamic programming analysis 
of optimal effort choice under differing productivity profiles

\section{Utility Maximization Framework}

To connect the empirical production estimates to an economic theory of
student effort, we model each student as choosing daily study hours $H$
to maximize a per-period payoff that values academic progress but
penalizes study effort. Let $G$ denote current GPA. Next-period GPA 
evolves according to
\[
G' = G + f(H; d),
\]
where $d \in \{0,1\}$ indexes depression status and $f(H;d)$ is the 
estimated GPA increment from studying $H$ hours. The function $f(H;d)$
differs across variants of the empirical specification, and the dynamic
model incorporates each of these forms.

Students obtain utility from higher GPA but incur a disutility from 
studying. In the implementation, the payoff from GPA is assumed to be 
\emph{concave}, reflecting diminishing marginal utility of academic 
achievement. Specifically, letting $u(\cdot)$ denote the GPA utility 
function, the instantaneous payoff is
\[
U(G,H;d) = u(G') \;-\; \chi_d H,
\]
where $\chi_d>0$ is the marginal disutility of effort for type $d$.
Consistent with the code, $u(\cdot)$ follows the CRRA family:
\[
u(G) =
\begin{cases}
\log(G), & \text{if } \sigma = 1,\\[6pt]
\dfrac{G^{1-\sigma}-1}{\,1-\sigma\,}, & \text{if } \sigma \neq 1,
\end{cases}
\]
with $\sigma=1$ (log utility) used as the baseline specification.

In dynamic form, students choose $H$ to maximize discounted lifetime
utility:
\[
V(G;d)
=
\max_{H \in [0,12]}
\left\{
u(G') - \chi_d H
\;+\;
\beta V(G';d)
\right\},
\qquad 0<\beta<1.
\]

The code implements this dynamic problem by discretizing the GPA state
space and solving for the optimal policy $H^*(G;d)$ using value
iteration. Because the empirical production function admits several
parametric forms, we evaluate optimal study policies under multiple
transition specifications $f(H;d)$.


\section{Dynamic Programming Model}
\label{sec:dp}

We embed the empirical productivity differences into a simple dynamic
programming (DP) setting. The state variable is GPA $G$ and the control
is study time $h$. The code computes four alternative versions of the
transition function, reflecting different assumptions about how GPA
responds to hours.

\subsection{Empirical Transition Structures}

Let $G_t$ denote GPA at time $t$. The model computes optimal study hours
under three specifications for the GPA increment $f(H;d)$:

\begin{enumerate}
    \item \textbf{Linearized productivity (LIN)}:
    \[
    f_{\mathrm{lin}}(H;d) = MP_d(\bar{H}) \cdot H,
    \]
    where $MP_d(\bar{H})$ is the empirical marginal product for type $d$
    evaluated at mean study hours.

    \item \textbf{Quadratic full (QF)}:
    \[
    f_{\mathrm{qf}}(H;d)
    =
    (b_1 + g_{2} d) H
    +
    (b_2 + g_{3} d) H^2,
    \]
    where $(b_1,b_2,g_2,g_3)$ are estimated from the empirical regression.


    \item \textbf{Centered regression (CC)}:
    \[
    f_{\mathrm{cc}}(H;d)
    =
    (b_1^c + g_2^c d) H_c
    +
    (b_2^c + g_3^c d) H_c^2,
    \qquad
    H_c = H - \bar{H},
    \]
    where coefficients $(b_1^c,b_2^c,g_2^c,g_3^c)$ come from the centered
    polynomial regression estimated in the code. This variant typically
    produces the most stable marginal product measures.
\end{enumerate}

Each specification yields a different state transition
\[
G_{t+1} = G_t + f_{\ell}(H_t;d),
\]
with $\ell \in \{\mathrm{lin},\mathrm{qf},\mathrm{cc}\}$.

\subsection{Instantaneous Utility}

In all specifications we treat effort cost additively but allow the
student's payoff from GPA to be \emph{concave}. Denote by $u(\cdot)$ the
per-period utility derived from next-period GPA $G_{t+1}$. The instantaneous
payoff is therefore
\[
u(G_{t+1}) - \chi_d h_t.
\]

In the code we implement the standard CRRA family for $u(\cdot)$, with
parameter $\sigma$:
\[
u(G) =
\begin{cases}
\log(G) & \text{if }\sigma = 1,\\[4pt]
\dfrac{G^{1-\sigma}-1}{1-\sigma} & \text{if }\sigma \neq 1,
\end{cases}
\]
and we impose a small numerical floor to ensure $G>0$ inside the utility
function. The linear disutility of effort remains $\chi_d h_t$, where
$\chi_d$ is calibrated from the empirical marginal product at mean
hours as in the code:
\[
\chi_0 = MP_{\mathrm{non}}(\bar{H}),\qquad
\chi_1 = MP_{\mathrm{dep}}(\bar{H}).
\]

In all numerical results presented below, we set $\sigma = 1$ so that
$u(G)=\log(G)$, following the instructor’s recommendation for a smooth
and concave specification.


\section{Policy Functions Across Specifications}
\label{sec:dp_results}

The DP solver produces 3 pairs of policy functions:
\[
\{h^*_{\mathrm{lin},d}(G), \;
  h^*_{\mathrm{qf},d}(G), \;
  h^*_{\mathrm{cc},d}(G)\}, \qquad d \in \{0,1\}.
\]

A key result from the code is that \emph{the relationship between
depression and optimal hours is specification-dependent}. In some cases,
such as the full quadratic model, depressed students choose lower hours
at certain GPA levels due to the curvature of the estimated production
function. In others---most notably the centered polynomial and linear
variants---the depressed policy lies above the non-depressed policy at
many GPA values.
Because the GPA payoff is concave in the implemented model (log/CRRA
utility), the resulting policy functions are smoother than under a
linear-utility formulation. Concavity reduces the prevalence of corner
solutions—situations in which students study either the maximum or
minimum feasible hours—and instead yields gradually declining optimal
hours as GPA rises. This aligns with economic intuition: when the
marginal utility of GPA diminishes at higher levels, additional study
effort becomes less valuable, leading to interior optima.


Evaluating the policies at the sample mean GPA 
($\overline{G}\approx 7.656$) yields, for example:

\[
h^*_{\mathrm{lin},0}(\overline{G}) \approx 8.75, \qquad
h^*_{\mathrm{lin},1}(\overline{G}) \approx 12.0,
\]

but the quadratic specifications may produce smaller or even reversed
differences. With concave utility, these reversals are less extreme than
under linear utility, but the qualitative ranking can still vary across
specifications because the productivity functions imply different
returns to effort across the GPA range. This reflects the fact that each
empirical specification implies a different marginal productivity
schedule.


Across all specifications, however, two qualitative observations hold:

\begin{itemize}
    \item Policies tend to be weakly decreasing in GPA: students with
          higher GPA need fewer additional study hours to sustain
          performance.
    \item Depression affects optimal hours through two competing
          channels: reduced productivity ($f_{\ell}$ smaller for $d=1$)
          and reduced effort cost ($\chi_1 < \chi_0$). The balance of
          these forces determines the ranking of $h^*(G;0)$ and
          $h^*(G;1)$, yielding specification-dependent results.
\end{itemize}

Figure 5 and 6 plots one representative set of the computed
policies.

% ---- main DP comparison (replace old dp_opt_policy)
\begin{figure}[H]
\centering
\includegraphics[width=0.85\textwidth]{./analysis_outputs/dp_policy_comparison_logutil.png}
\caption{Comparison of optimal study-hour policies under three empirical transition specifications (linearized, quadratic full, centered). Lines compare non-depressed and depressed students; the ranking of effort is specification-dependent.}
\label{fig:dp_policy_comparison}
\end{figure}

% ---- value functions (centered spec)
\begin{figure}[H]
\centering
\includegraphics[width=0.8\textwidth]{./analysis_outputs/dp_value_centered_logutil.png}
\caption{Value functions under the centered transition specification with log utility. The concave utility causes smooth value curves and reduces corner solutions.}
\label{fig:dp_value_centered}
\end{figure}



\section{Conclusion}

The dynamic programming analysis demonstrates how different empirical 
assumptions about academic productivity translate into different optimal
study behaviors. Because depression reduces estimated marginal
productivity, optimal study hours may increase or decrease depending on
how strongly the cost parameter $\chi_d$ adjusts in the calibration and
how the production function is modeled.

The code's multiple specifications show that predictions are 
\emph{sensitive to the chosen functional form}. In particular, while some
specifications predict higher optimal hours for depressed students, 
others predict smaller or reversed gaps. This highlights the importance 
of carefully modeling the productivity–effort relationship when using 
dynamic models to interpret behavioral responses to mental health 
conditions.





\end{document}